\rtmtight
\begin{tblr}{width=\textwidth,colspec={|Q[c,m,2.1cm]|X[l,m]|},hlines,vlines,hspan=minimal,rowsep=1.5pt,colsep=3pt,row{1}={bg=headgray,font=\bfseries}}
\SetCell{c,m}\hfil\logo\hfil & \textbf{POLITEKNIK NEGERI MALANG}\par\textbf{JURUSAN TEKNOLOGI INFORMASI}\par\textbf{PROGRAM STUDI: D4 TEKNIK INFORMATIKA} \\
\SetCell[c=2]{l} \textbf{RENCANA TUGAS MAHASISWA (RTM) DAN RUBRIK PENILAIAN} & \\
Nama Mata Kuliah & Rekayasa Sistem \\
Kode Mata Kuliah & RTI256206 \quad \textbf{sks / Jam perminggu:} 6 SKS/12 jam \quad \textbf{Semester:} 6 \\
Dosen Pengampu & Tim Pengajar Program Studi D4 TI \\
\end{tblr}
\assessmentblock{Dokumen Arsitektur Sistem}{Case Method (CM) -- perancangan arsitektur terstruktur}{SCPMK0211-05501, SCPMK0602-05502}{Mahasiswa menyusun dokumen arsitektur sistem enterprise yang komprehensif, mencakup diagram komponen, antarmuka, aliran data, keputusan arsitektur (ADR), dan spesifikasi kebutuhan awal.}{Menganalisis konteks sistem, merancang komponen dan antarmuka, mendokumentasikan keputusan arsitektur, dan menyusun kebutuhan awal dalam format SRS.}{Dokumen arsitektur (C4 diagram/UML), ADR, spesifikasi kebutuhan awal, dan presentasi arsitektur.}{Kelengkapan dan kualitas desain arsitektur & 6 & Komponen, antarmuka, aliran data, dan ketergantungan antar modul dirancang dengan benar dan terdokumentasi dalam diagram yang dapat dibaca. & Desain arsitektur mencakup komponen utama, tetapi beberapa antarmuka atau aliran data belum terdokumentasi lengkap. & Desain arsitektur tidak lengkap, diagram tidak dapat dibaca, atau tidak mencerminkan kebutuhan sistem. \\ \\
Kualitas kebutuhan awal dan traceability & 5 & Kebutuhan fungsional dan non-fungsional terdefinisi jelas, dapat diverifikasi, dan terhubung ke komponen arsitektur. & Kebutuhan sebagian besar jelas tetapi traceability ke arsitektur atau verifikabilitas masih belum lengkap. & Kebutuhan tidak terdefinisi jelas, tidak dapat diverifikasi, atau tidak terhubung ke arsitektur. \\ \\
Konsistensi dokumentasi dan keterbacaan & 4 & Dokumen konsisten, mengikuti standar notasi, menggunakan template yang seragam, dan dapat dipahami oleh semua stakeholder. & Dokumentasi cukup konsisten tetapi beberapa bagian belum mengikuti standar notasi atau masih sulit dipahami. & Dokumentasi tidak konsisten, tidak mengikuti standar, atau tidak dapat dipahami stakeholder. \\}

\assessmentblock{Sprint Rekayasa Sistem}{Project Based Learning (PBL) -- Sprint 1-2, UTS Milestone, Sprint 3-4}{SCPMK0211-05501, SCPMK0602-05502, SCPMK0609-05503}{Mahasiswa mengeksekusi sprint rekayasa sistem meliputi validasi arsitektur, implementasi pipeline CI/CD, strategi deployment, dan UTS milestone berupa presentasi arsitektur dan SRS kepada stakeholder.}{Melaksanakan sprint planning, memvalidasi arsitektur dengan stakeholder, mengimplementasikan CI/CD, merancang strategi deployment, dan mendokumentasikan setiap sprint.}{Artefak arsitektur tervalidasi, pipeline CI/CD berjalan, laporan sprint, bukti demo, dan evidence UTS milestone.}{Desain arsitektur dan validasi komponen & 18 & Arsitektur divalidasi bersama stakeholder, komponen diimplementasikan sesuai desain, dan aliran data terkonfirmasi berfungsi. & Validasi arsitektur dilakukan tetapi beberapa komponen atau aliran data belum sepenuhnya dikonfirmasi. & Arsitektur tidak divalidasi, komponen tidak sesuai desain, atau aliran data tidak berfungsi. \\ \\
Implementasi pipeline CI/CD dan automated testing & 20 & Pipeline CI/CD berjalan end-to-end, automated test suite mencakup unit dan integration test, dan build gagal ketika test tidak lulus. & Pipeline CI/CD berjalan untuk alur utama tetapi automated test masih terbatas atau belum terintegrasi penuh. & Pipeline CI/CD tidak berjalan, automated testing tidak ada, atau build tidak memberikan feedback yang bermakna. \\ \\
Strategi deployment dan operasi sistem & 17 & Strategi blue-green/canary diimplementasikan, rollback terdokumentasi, dan monitoring/observability terkonfigurasi. & Strategi deployment ada tetapi rollback atau monitoring masih belum lengkap atau teruji. & Tidak ada strategi deployment yang terstruktur atau sistem tidak dapat di-deploy dengan aman. \\}

\assessmentblock{Laporan SRS}{Case Method (CM) -- dokumen spesifikasi kebutuhan}{SCPMK0602-05502}{Mahasiswa menyusun dokumen SRS (Software Requirements Specification) final yang telah divalidasi stakeholder, mencakup kebutuhan fungsional, non-fungsional, batasan sistem, dan acceptance criteria.}{Mengkonsolidasikan hasil elicitation, menyusun SRS sesuai standar IEEE 830, melakukan review bersama stakeholder, dan mendokumentasikan hasil sign-off.}{Dokumen SRS final sesuai template standar, hasil review stakeholder, acceptance criteria, dan bukti sign-off.}{Kelengkapan dan kualitas dokumen SRS & 4 & SRS mencakup semua kebutuhan fungsional dan non-fungsional, batasan sistem terdefinisi, dan setiap kebutuhan dapat diverifikasi. & SRS mencakup kebutuhan utama tetapi beberapa kebutuhan non-fungsional atau batasan masih belum lengkap. & SRS tidak mencakup kebutuhan penting atau banyak kebutuhan yang tidak dapat diverifikasi. \\ \\
Validasi dan sign-off stakeholder & 3 & Stakeholder mereview dan menandatangani SRS, komentar ditindaklanjuti, dan versi final terdistribusi ke semua pihak. & Review stakeholder dilakukan tetapi beberapa komentar belum ditindaklanjuti atau sign-off belum lengkap. & Tidak ada proses review stakeholder atau SRS tidak mendapat persetujuan yang diperlukan. \\ \\
Traceability kebutuhan ke arsitektur & 3 & Setiap kebutuhan dalam SRS dapat ditelusuri ke komponen arsitektur dan test case yang mengujinya. & Traceability tersedia untuk kebutuhan utama tetapi tidak semua kebutuhan terhubung ke arsitektur atau test case. & Tidak ada traceability antara kebutuhan, arsitektur, dan pengujian. \\}

\assessmentblock{UAS Defense Rekayasa Sistem}{UAS defense arsitektur dan demo deployment}{SCPMK0211-05501, SCPMK0609-05503}{Mahasiswa melakukan defense arsitektur sistem enterprise yang dirancang dan mendemonstrasikan pipeline deployment yang diimplementasikan kepada panel penguji akademik dan industri.}{Mempersiapkan materi defense, mendemonstrasikan arsitektur sistem dan pipeline deployment, serta menjawab pertanyaan teknis mengenai keputusan desain dan implementasi.}{Materi presentasi, demonstrasi arsitektur dan deployment pipeline yang berjalan, dan dokumentasi sistem akhir.}{Kualitas arsitektur dan demonstrasi deployment & 8 & Arsitektur enterprise terdokumentasi lengkap, pipeline deployment berjalan end-to-end, dan demo menunjukkan sistem berfungsi di lingkungan produksi. & Arsitektur dan pipeline deployment ada tetapi beberapa aspek belum terdemonstrasikan sepenuhnya. & Arsitektur tidak terdokumentasi lengkap atau pipeline deployment tidak dapat didemonstrasikan. \\ \\
Kemampuan defense keputusan teknis & 7 & Keputusan arsitektur dan deployment dijelaskan dengan argumentasi berbasis trade-off analysis, best practice, dan kebutuhan sistem. & Defense cukup jelas tetapi beberapa keputusan teknis belum didukung argumentasi yang kuat. & Tidak mampu menjelaskan keputusan teknis atau argumentasi tidak berdasarkan prinsip rekayasa. \\ \\
Kelengkapan dokumentasi sistem akhir & 5 & Dokumentasi sistem akhir mencakup ADR, diagram arsitektur terkini, runbook deployment, dan laporan operasi sistem. & Dokumentasi sistem tersedia tetapi beberapa komponen seperti runbook atau ADR masih belum lengkap. & Dokumentasi sistem tidak lengkap atau tidak mencukupi untuk memelihara sistem setelah delivery. \\}

\begin{tblr}{width=\textwidth,colspec={|Q[l,m,3.2cm]|X[l,m]|},hlines,vlines,hspan=minimal,row{1-3}={bg=headgray},rowsep=1.5pt,colsep=3pt}
Jadwal Pelaksanaan: & Minggu 1--17 sesuai RPS. Setiap evaluasi memiliki RTM dan rubrik multi-kriteria. \\
Lain-Lain Yang Diperlukan: & LMS, repositori kode sumber, tools CI/CD dan diagramming, perangkat lunak sesuai mata kuliah, dan perangkat presentasi. \\
Pustaka: & Mengacu pustaka utama dan pendukung pada bagian RPS. \\
\end{tblr}
